\begin{theoremEnd}[combined main,category=joint]{proposition}[Joint sign constraints]
\CombinedResultCheck{prop:joint}{proposition}\label{prop:joint}
Under the finite Gaussian model in \cref{sec:finite-model}, let
$\boldsymbol a=(a_0,\ldots,a_m)^T$ denote the coefficient vector and fix
a nonempty $A\subseteq\{0,1,2\}$ and signs
$\varepsilon:A\to\{-1,1\}$, with event $\mathcal C_A^\varepsilon$ defined in
\cref{eq:joint-event}. Suppose an invertible linear change of coordinates
$\boldsymbol a\mapsto(\boldsymbol g,\gamma)\in\R^m\times\R$
gives $f_{\boldsymbol a}=f_{\boldsymbol g}+\gamma h$.
Here $h$ is a fixed function in the retained span, shared by all selected
orders, $\gamma$ is its scalar coefficient, and $f_{\boldsymbol g}$
depends linearly on the remaining coordinates $\boldsymbol g$.
Assume $h,f_{\boldsymbol g}\in C^{\max A}(I)$ with
$\min_{x\in I}\varepsilon_r\frac{\mathrm{d}^rh}{\mathrm{d}x^r}(x)>0$
for every $r\in A$. Define
\begin{equation}\label{eq:joint-threshold}
 \eta_r(x;\boldsymbol g)=
       \frac{-\varepsilon_r\frac{\mathrm{d}^rf_{\boldsymbol g}}{\mathrm{d}x^r}(x)}
            {\varepsilon_r\frac{\mathrm{d}^rh}{\mathrm{d}x^r}(x)},
 \qquad
 \tau_A^\varepsilon(\boldsymbol g)=
       \max_{r\in A}\max_{x\in I}\eta_r(x;\boldsymbol g).
\end{equation}
This finite, attained threshold satisfies
$f_{\boldsymbol a}\in\mathcal C_A^\varepsilon$ exactly when
$\gamma\geq\tau_A^\varepsilon(\boldsymbol g)$.
The event has positive ordinary posterior probability and a null
coefficient boundary. Its marginal density and reconstruction are
\cref{eq:marginal,eq:reconstruction}, with $\tau_r,\mathcal C_r$
replaced by $\tau_A^\varepsilon,\mathcal C_A^\varepsilon$ and the
Gaussian parameters in \cref{eq:conditional} computed under this map.
The conditioned law satisfies
$\min_{x\in I}\varepsilon_r\frac{\mathrm{d}^rf_{\boldsymbol a}}{\mathrm{d}x^r}(x)>0$
almost surely for every $r\in A$. The envelope, strong log-concavity,
and conservative threshold-error conclusions extend under their
respective hypotheses.
\end{theoremEnd}

\begin{proofEnd}
\CombinedProofCheck{prop:joint}
Fix the coefficient map and signed direction satisfying the hypotheses
of \cref{prop:joint}.
Dividing each signed derivative inequality by its positive direction
derivative gives $\gamma\geq\eta_r(x;\boldsymbol g)$, so their intersection
gives \cref{eq:joint-threshold}. Compactness of $I$, finiteness
of $A$, and the denominator bounds make the maximum finite and attained.
The Gaussian conditioning argument in the proof of \cref{thm:finite} then gives
positive probability, the marginal density, and reconstruction.

Above the threshold, every selected signed derivative has a strictly
positive minimum. At equality, an attained maximizing pair gives a zero
derivative, and perturbing $\gamma$ crosses the event boundary. Thus the
boundary is the threshold graph, which has zero conditional Gaussian
probability for every $\boldsymbol g$, establishing almost-sure strictness. Each
$\eta_r(x;\boldsymbol g)$ is linear in $\boldsymbol g$, so the envelope
and strong log-concavity arguments apply with the compact index set
$A\times I$, and \cref{prop:error} applies to the same scalar truncation.
\end{proofEnd}
